\begin{frame}{그림으로: 힌지 손실의 모양}
  \begin{center}
  \includegraphics[height=5.3cm]{ch05_2_hinge_loss.png}
  \end{center}
  \vskip0.2em
  \footnotesize
  \begin{itemize}
    \item $f\ge1$(마진 밖)이면 손실이 \textbf{정확히 0} --- 힌지의 핵심 성질.
    \item $f=1$에서 기울기가 꺾이는 지점(kink)이 ``힌지''라는 이름의 유래.
    \item 점선(교차 엔트로피)은 아무리 잘 맞혀도 손실이 0이 안 됨 --- 이
      차이가 서포트 벡터 개념을 가능하게 한다.
  \end{itemize}
\end{frame}
